\documentclass[a4paper,12pt,twoside]{memoir}

% Castellano
\usepackage[spanish,es-tabla,es-nodecimaldot]{babel}
\selectlanguage{spanish}
\usepackage[utf8]{inputenc}
\usepackage[T1]{fontenc}
\usepackage{lmodern} % scalable font
\usepackage{microtype}
\usepackage{placeins}

\RequirePackage{booktabs}
\RequirePackage[table]{xcolor}
\RequirePackage{xtab}
\RequirePackage{multirow}

\usepackage{dirtree}
\renewcommand\DTstylecomment{\rmfamily}

% Links
\PassOptionsToPackage{hyphens}{url}\usepackage[colorlinks]{hyperref}
\hypersetup{
	allcolors = {red}
}

% Ecuaciones
\usepackage{amsmath}

% Rutas de fichero / paquete
\newcommand{\ruta}[1]{{\sffamily #1}}

% Párrafos
\nonzeroparskip

% Huérfanas y viudas
\widowpenalty100000
\clubpenalty100000

% Evitar solapes en el header
\nouppercaseheads

% Imagenes
\usepackage{graphicx}
\usepackage{adjustbox}
\newcommand{\imagen}[3]{
	\begin{figure}[!h]
		\centering
		\includegraphics[width=#3\textwidth]{#1}
		\caption{#2}\label{fig:#1}
	\end{figure}
	\FloatBarrier
}

\newcommand{\imagenflotante}[2]{
	\begin{figure}%[!h]
		\centering
		\includegraphics[width=0.9\textwidth]{#1}
		\caption{#2}\label{fig:#1}
	\end{figure}
}



% El comando \figura nos permite insertar figuras comodamente, y utilizando
% siempre el mismo formato. Los parametros son:
% 1 -> Porcentaje del ancho de página que ocupará la figura (de 0 a 1)
% 2 --> Fichero de la imagen
% 3 --> Texto a pie de imagen
% 4 --> Etiqueta (label) para referencias
% 5 --> Opciones que queramos pasarle al \includegraphics
% 6 --> Opciones de posicionamiento a pasarle a \begin{figure}
\newcommand{\figuraConPosicion}[6]{%
  \setlength{\anchoFloat}{#1\textwidth}%
  \addtolength{\anchoFloat}{-4\fboxsep}%
  \setlength{\anchoFigura}{\anchoFloat}%
  \begin{figure}[#6]
    \begin{center}%
      \Ovalbox{%
        \begin{minipage}{\anchoFloat}%
          \begin{center}%
            \includegraphics[width=\anchoFigura,#5]{#2}%
            \caption{#3}%
            \label{#4}%
          \end{center}%
        \end{minipage}
      }%
    \end{center}%
  \end{figure}%
}

%
% Comando para incluir imágenes en formato apaisado (sin marco).
\newcommand{\figuraApaisadaSinMarco}[5]{%
  \begin{figure}%
    \begin{center}%
    \includegraphics[angle=90,height=#1\textheight,#5]{#2}%
    \caption{#3}%
    \label{#4}%
    \end{center}%
  \end{figure}%
}
% Para las tablas
\newcommand{\otoprule}{\midrule [\heavyrulewidth]}
%
% Nuevo comando para tablas pequeñas (menos de una página).
\newcommand{\tablaSmall}[5]{%
 \begin{table}[h!]
  \begin{center}
   \rowcolors {2}{gray!35}{}
   \begin{tabular}{#2}
    \toprule
    #4
    \otoprule
    #5
    \bottomrule
   \end{tabular}
   \caption{#1}
   \label{tabla:#3}
  \end{center}
 \end{table}
}

%
%Para el float H de tablaSmallSinColores
\usepackage{float}

%
% Nuevo comando para tablas pequeñas (menos de una página).
\newcommand{\tablaSmallSinColores}[5]{%
 \begin{table}[H]
  \begin{center}
   \begin{tabular}{#2}
    \toprule
    #4
    \otoprule
    #5
    \bottomrule
   \end{tabular}
   \caption{#1}
   \label{tabla:#3}
  \end{center}
 \end{table}
}

\newcommand{\tablaApaisadaSmall}[5]{%
\begin{landscape}
  \begin{table}
   \begin{center}
    \rowcolors {2}{gray!35}{}
    \begin{tabular}{#2}
     \toprule
     #4
     \otoprule
     #5
     \bottomrule
    \end{tabular}
    \caption{#1}
    \label{tabla:#3}
   \end{center}
  \end{table}
\end{landscape}
}

%
% Nuevo comando para tablas grandes con cabecera y filas alternas coloreadas en gris.
\newcommand{\tabla}[6]{%
  \begin{center}
    \tablefirsthead{
      \toprule
      #5
      \otoprule
    }
    \tablehead{
      \multicolumn{#3}{l}{\small\sl continúa desde la página anterior}\\
      \toprule
      #5
      \otoprule
    }
    \tabletail{
      \hline
      \multicolumn{#3}{r}{\small\sl continúa en la página siguiente}\\
    }
    \tablelasttail{
      \hline
    }
    \bottomcaption{#1}
    \rowcolors {2}{gray!35}{}
    \begin{xtabular}{#2}
      #6
      \bottomrule
    \end{xtabular}
    \label{tabla:#4}
  \end{center}
}

%
% Nuevo comando para tablas grandes con cabecera.
\newcommand{\tablaSinColores}[6]{%
  \begin{center}
    \tablefirsthead{
      \toprule
      #5
      \otoprule
    }
    \tablehead{
      \multicolumn{#3}{l}{\small\sl continúa desde la página anterior}\\
      \toprule
      #5
      \otoprule
    }
    \tabletail{
      \hline
      \multicolumn{#3}{r}{\small\sl continúa en la página siguiente}\\
    }
    \tablelasttail{
      \hline
    }
    \bottomcaption{#1}
    \begin{xtabular}{#2}
      #6
      \bottomrule
    \end{xtabular}
    \label{tabla:#4}
  \end{center}
}

%
% Nuevo comando para tablas grandes sin cabecera.
\newcommand{\tablaSinCabecera}[5]{%
  \begin{center}
    \tablefirsthead{
      \toprule
    }
    \tablehead{
      \multicolumn{#3}{l}{\small\sl continúa desde la página anterior}\\
      \hline
    }
    \tabletail{
      \hline
      \multicolumn{#3}{r}{\small\sl continúa en la página siguiente}\\
    }
    \tablelasttail{
      \hline
    }
    \bottomcaption{#1}
  \begin{xtabular}{#2}
    #5
   \bottomrule
  \end{xtabular}
  \label{tabla:#4}
  \end{center}
}



\definecolor{cgoLight}{HTML}{EEEEEE}
\definecolor{cgoExtralight}{HTML}{FFFFFF}

%
% Nuevo comando para tablas grandes sin cabecera.
\newcommand{\tablaSinCabeceraConBandas}[5]{%
  \begin{center}
    \tablefirsthead{
      \toprule
    }
    \tablehead{
      \multicolumn{#3}{l}{\small\sl continúa desde la página anterior}\\
      \hline
    }
    \tabletail{
      \hline
      \multicolumn{#3}{r}{\small\sl continúa en la página siguiente}\\
    }
    \tablelasttail{
      \hline
    }
    \bottomcaption{#1}
    \rowcolors[]{1}{cgoExtralight}{cgoLight}

  \begin{xtabular}{#2}
    #5
   \bottomrule
  \end{xtabular}
  \label{tabla:#4}
  \end{center}
}




\graphicspath{ {./img/} }

% Capítulos
\chapterstyle{bianchi}
\newcommand{\capitulo}[2]{
	\setcounter{chapter}{#1}
	\setcounter{section}{0}
	\setcounter{figure}{0}
	\setcounter{table}{0}
	\chapter*{#2}
	\addcontentsline{toc}{chapter}{#2}
	\markboth{#2}{#2}
}

% Apéndices
\renewcommand{\appendixname}{Apéndice}
\renewcommand*\cftappendixname{\appendixname}

\newcommand{\apendice}[1]{
	%\renewcommand{\thechapter}{A}
	\chapter{#1}
}

\renewcommand*\cftappendixname{\appendixname\ }

% Formato de portada
\makeatletter
\usepackage{xcolor}
\newcommand{\tutor}[1]{\def\@tutor{#1}}
\newcommand{\course}[1]{\def\@course{#1}}
\definecolor{cpardoBox}{HTML}{E6E6FF}
\def\maketitle{
  \null
  \thispagestyle{empty}
  % Cabecera ----------------
\noindent\includegraphics[width=\textwidth]{cabecera}\vspace{1cm}%
  \vfill
  % Título proyecto y escudo informática ----------------
  \colorbox{cpardoBox}{%
    \begin{minipage}{.8\textwidth}
      \vspace{.5cm}\Large
      \begin{center}
      \textbf{TFG del Grado en Ingeniería Informática}\vspace{.6cm}\\
      \textbf{\LARGE\@title{}}
      \end{center}
      \vspace{.2cm}
    \end{minipage}

  }%
  \hfill\begin{minipage}{.20\textwidth}
    \includegraphics[width=\textwidth]{escudoInfor}
  \end{minipage}
  \vfill
  % Datos de alumno, curso y tutores ------------------
  \begin{center}%
  {%
    \noindent\LARGE
    Presentado por \@author{}\\ 
    en Universidad de Burgos --- \@date{}\\
    Tutor: \@tutor{}\\
  }%
  \end{center}%
  \null
  \cleardoublepage
  }
\makeatother


% Datos de portada
\title{Detección de \textit{malware} mediante técnicas de Inteligencia Artificial \\Documentación Técnica}
\author{David Cezar Toderas}
\tutor{Álvar Arnaiz González}
\date{\today}

\begin{document}

\maketitle



\cleardoublepage



%%%%%%%%%%%%%%%%%%%%%%%%%%%%%%%%%%%%%%%%%%%%%%%%%%%%%%%%%%%%%%%%%%%%%%%%%%%%%%%%%%%%%%%%



\frontmatter


\clearpage

% Indices
\tableofcontents

\clearpage

\listoffigures

\clearpage

\listoftables

\clearpage

\mainmatter

\appendix

\apendice{Plan de Proyecto Software}

\section{Introducción}

\section{Planificación temporal}

\section{Estudio de viabilidad}

\subsection{Viabilidad económica}

\subsection{Viabilidad legal}



\apendice{Especificación de Requisitos}

\section{Introducción}

Este apéndice documenta la especificación de requisitos del trabajo desarrollado. Su propósito es definir de manera formal y detallada las capacidades, características y restricciones del sistema. Se establecen los objetivos generales del proyecto, se desglosa un catálogo de requisitos funcionales y no funcionales, y finalmente, se describen los principales casos de uso que ilustran la interacción del usuario con la aplicación.

\section{Objetivos generales}
Los objetivos generales representan las metas de alto nivel que se buscaron alcanzar con la realización de este trabajo, combinando tanto las aspiraciones de investigación como los entregables prácticos.

\begin{enumerate}
	\item \textbf{Investigar el estado del arte:} Realizar un análisis de la literatura científica para comprender las técnicas actuales de detección de \textit{malware} con IA y posicionar el proyecto en el panorama actual.
	
	\item \textbf{Desarrollar un sistema de detección de extremo a extremo:} Construir un \textit{pipeline} completo, desde la recolección de datos y la extracción de características hasta el entrenamiento y la evaluación de un modelo funcional.
	
	\item \textbf{Alcanzar un alto rendimiento predictivo:} Lograr que los modelos desarrollados obtengan métricas de clasificación altas, con un enfoque especial maximizar el \textit{recall} para minimizar los falsos negativos.
	
	\item \textbf{Comparar diferentes arquitecturas de modelos:} Evaluar y contrastar el rendimiento de una red neuronal profunda frente a algoritmos de aprendizaje automático clásicos para obtener conclusiones sobre la eficacia de cada enfoque.
	
	\item \textbf{Garantizar la interpretabilidad del sistema:} Implementar técnicas que permitan explicar las decisiones de los modelos, aportando transparencia y confianza a los resultados.
	
	\item \textbf{Crear una aplicación de demostración:} Desarrollar una interfaz web interactiva que permita a un usuario probar y visualizar el funcionamiento de todo el sistema.
\end{enumerate}

\section{Catálogo de requisitos}
A continuación se presenta el catálogo detallado de requisitos que el sistema debe satisfacer.

\subsection{Requisitos funcionales}
Los requisitos funcionales (RF) especifican lo que el sistema debe hacer. Describen las funcionalidades, tareas y servicios que la aplicación final debe proporcionar al usuario para cumplir con su propósito.

\begin{itemize}
	\item \textbf{RF-1: Carga de archivos APK.} El sistema debe proporcionar una interfaz que permita al usuario seleccionar y subir un archivo con formato \texttt{.apk} para su posterior análisis.
	
	\item \textbf{RF-2: Análisis y clasificación de la aplicación.} Una vez subida una APK, el sistema debe ejecutar el \textit{pipeline} de análisis completo. Esto incluye la extracción de características estáticas, el preprocesamiento de los datos a través del \textit{embedder} y la ejecución de la inferencia con todos los modelos entrenados (la red neuronal y los clasificadores clásicos).
	
	\item \textbf{RF-3: Visualización de las predicciones.} La aplicación debe presentar al usuario los resultados de la clasificación de forma clara. Debe mostrar un veredicto general y una tabla detallada con la predicción de cada modelo individual, incluyendo los porcentajes de confianza para las clases <<benigno>> y <<malicioso>>.
	
	\begin{itemize}
		\item \textbf{RF-3.1: Transparencia del proceso de extracción.} Para que el proceso no sea una <<caja negra>>, la interfaz debe permitir al usuario inspeccionar las características <<brutas>> que han sido extraídas de la APK. Esto incluye las diferentes listas de permisos, actividades, servicios, receptores, así como propiedades del archivo como su tamaño o su \textit{fuzzy hash}.
		
		\item \textbf{RF-3.2: Visualización de datos procesados.} Además de los datos en <<bruto>>, el sistema debe mostrar la representación numérica en la que se transforman. Esto implica visualizar tanto los datos \textit{tokenizados} y escalados como el vector de \textit{embeddings} final que se introduce en los clasificadores.
	\end{itemize}
	
	\item \textbf{RF-4: Interpretabilidad de las predicciones.} La aplicación debe ofrecer explicaciones sobre las decisiones del modelo. Para ello, deberá generar y mostrar un conjunto de gráficos de SHAP, incluyendo la importancia global de las características y análisis locales que detallen qué factores han influido en la predicción de la muestra actual.
	
	\item \textbf{RF-5: Visualización del espacio de características.} El sistema debe ser capaz de generar una proyección 2D del espacio de \textit{embeddings} mediante UMAP. En esta visualización, se debe mostrar la distribución de las muestras del \textit{dataset} de fondo y resaltar la posición de la APK recién analizada, permitiendo al usuario entender su ubicación relativa respecto a las clases conocidas.
	
	\item \textbf{RF-6: Gestión del historial de análisis.} El sistema debe mantener un historial de las aplicaciones analizadas durante la sesión activa del usuario. La interfaz debe permitir al usuario seleccionar una entrada del historial para volver a cargar sus resultados y ofrecer una opción para borrar todo el historial de la sesión.
\end{itemize}

\subsection{Requisitos no funcionales}
Los requisitos no funcionales (RNF) describen los atributos de calidad y las restricciones bajo las cuales el sistema debe operar. No se refieren a qué hace el sistema, sino a <<cómo>> lo hace, definiendo aspectos como su rendimiento, fiabilidad o portabilidad.

\begin{itemize}
	\item \textbf{RNF-1: Rendimiento y fiabilidad de los modelos.} Los distintos clasificadores deben alcanzar un alto nivel de rendimiento, definido por métricas de evaluación estándar. Específicamente, se establece como requisito clave obtener una métrica de \textit{recall} cercana al 98\%, garantizando una detección muy alta de las muestras maliciosas.
	
	\item \textbf{RNF-2: Eficiencia del análisis.} El tiempo total desde que el usuario sube una APK hasta que recibe un resultado completo (incluyendo la extracción, el preprocesamiento y la inferencia) debe ser razonable para una buena experiencia de usuario, idealmente completándose en menos de un par de minutos.
	
	\item \textbf{RNF-3: Portabilidad y facilidad de despliegue.} Todo el sistema, incluyendo la aplicación web, el modelo y sus dependencias, debe estar empaquetado en un contenedor Docker. Esto asegura que la aplicación sea portable y pueda ser desplegada de forma sencilla y consistente en diferentes sistemas operativos.
	
	\item \textbf{RNF-4: Usabilidad de la interfaz.} La aplicación web debe tener una interfaz de usuario clara, intuitiva y fácil de navegar para un perfil de usuario con conocimientos técnicos, pero no necesariamente experto en inteligencia artificial. La información debe presentarse de forma organizada y comprensible.
	
	\item \textbf{RNF-5: Modularidad del código.} La base de código del proyecto debe seguir un diseño modular que separe claramente las distintas responsabilidades (extracción de datos, arquitectura del modelo, entrenamiento, etc.). Esto facilita el mantenimiento, la experimentación y el posible desarrollo futuro del sistema.
\end{itemize}

\section{Casos de uso}
Los casos de uso describen las interacciones entre un actor (en este caso, el <<Usuario/Analista>>) y el sistema para alcanzar un objetivo. A continuación se presenta un diagrama esquemático (Figura \ref{fig:use_cases_diagram}) y se detallan los tres casos de uso principales.

% Diagrama Esquelético de Casos de Uso:
% Actor: Usuario/Analista
% Casos de Uso:
% 1. (Usuario/Analista) -> Analizar una aplicación APK
% 2. (Usuario/Analista) -> Interpretar resultados de una predicción
% 3. (Usuario/Analista) -> Gestionar historial de análisis

\begin{table}[H]
	\centering
	\begin{tabularx}{\linewidth}{ p{0.21\columnwidth} p{0.71\columnwidth} }
		\toprule
		\textbf{CU-1}    & \textbf{Analizar una nueva aplicación APK}\\
		\toprule
		\textbf{Versión}              & 1.0    \\
		\textbf{Autor}                & David Cezar Toderas \\
		\textbf{Requisitos asociados} & RF-1, RF-2, RF-3 \\
		\textbf{Descripción}          & El usuario sube un archivo APK al sistema para que este sea analizado por los modelos de IA y se muestren los resultados de la clasificación. \\
		\textbf{Precondición}         & El usuario tiene un archivo \texttt{.apk} válido y la aplicación web está en ejecución. \\
		\textbf{Acciones}             &
		\begin{enumerate}
			\def\labelenumi{\arabic{enumi}.}
			\tightlist
			\item El usuario accede a la aplicación web.
			\item El usuario arrastra y suelta un archivo APK en la zona de carga o lo selecciona mediante el explorador de archivos.
			\item El usuario presiona el botón <<Analizar APK>>.
			\item El sistema procesa el archivo: extrae las características, las pasa por el \textit{embedder} y realiza la inferencia con todos los modelos.
			\item La interfaz se actualiza para mostrar la pestaña de <<Predicciones>> con el veredicto de cada modelo.
			\item El análisis se añade al historial de la sesión.
		\end{enumerate}\\
		\textbf{Postcondición}        & Se muestra una predicción de clasificación para la APK subida. \\
		\textbf{Excepciones}          & \textbf{E-1.1:} Si el archivo subido no es un APK válido, el sistema muestra un mensaje de error. \newline \textbf{E-1.2:} Si el archivo supera el límite de tamaño, el sistema muestra una advertencia y no permite el análisis. \\
		\textbf{Importancia}          & Alta \\
		\bottomrule
	\end{tabularx}
	\caption{CU-1 Analizar una nueva aplicación APK.}
\end{table}


\begin{table}[H]
	\centering
	\begin{tabularx}{\linewidth}{ p{0.21\columnwidth} p{0.71\columnwidth} }
		\toprule
		\textbf{CU-2}    & \textbf{Interpretar el resultado de un análisis}\\
		\toprule
		\textbf{Versión}              & 1.0    \\
		\textbf{Autor}                & David Cezar Toderas \\
		\textbf{Requisitos asociados} & RF-3.1, RF-3.2, RF-4, RF-5 \\
		\textbf{Descripción}          & Tras analizar un APK, el usuario explora las diferentes pestañas de la interfaz para comprender en profundidad el resultado y el razonamiento del modelo. \\
		\textbf{Precondición}         & Se ha completado con éxito un análisis (CU-1). \\
		\textbf{Acciones}             &
		\begin{enumerate}
			\def\labelenumi{\arabic{enumi}.}
			\tightlist
			\item El usuario navega a la pestaña <<Características Extraídas>> para ver los datos brutos y procesados.
			\item El usuario navega a la pestaña <<Explicaciones>> para visualizar los gráficos de importancia de características de SHAP.
			\item El usuario navega a la pestaña <<UMAP>> para ver la proyección del \textit{embedding} de la aplicación analizada.
		\end{enumerate}\\
		\textbf{Postcondición}        & El usuario obtiene una visión detallada de los datos y las justificaciones que respaldan la predicción del modelo. \\
		\textbf{Excepciones}          & Ninguna. \\
		\textbf{Importancia}          & Alta \\
		\bottomrule
	\end{tabularx}
	\caption{CU-2 Interpretar el resultado de un análisis.}
\end{table}

\begin{table}[H]
	\centering
	\begin{tabularx}{\linewidth}{ p{0.21\columnwidth} p{0.71\columnwidth} }
		\toprule
		\textbf{CU-3}    & \textbf{Gestionar el historial de análisis}\\
		\toprule
		\textbf{Versión}              & 1.0    \\
		\textbf{Autor}                & David Cezar Toderas \\
		\textbf{Requisitos asociados} & RF-6 \\
		\textbf{Descripción}          & El usuario interactúa con el historial de análisis de la sesión actual para consultar resultados anteriores o para limpiar la lista. \\
		\textbf{Precondición}         & Se ha analizado al menos una aplicación durante la sesión actual. \\
		\textbf{Acciones}             &
		\textbf{Flujo A: Consultar un análisis anterior}
		\begin{enumerate}
			\def\labelenumi{\arabic{enumi}.}
			\tightlist
			\item El usuario hace clic sobre el nombre de un archivo en la lista del <<Historial de Análisis>>.
			\item El sistema carga en la vista principal todos los datos y resultados correspondientes a ese análisis.
		\end{enumerate}
		\textbf{Flujo B: Limpiar el historial}
		\begin{enumerate}
			\def\labelenumi{\arabic{enumi}.}
			\tightlist
			\item El usuario hace clic en el botón <<Limpiar Historial>>.
			\item El sistema borra todas las entradas del historial de la sesión.
		\end{enumerate}\\
		\textbf{Postcondición}        & \textbf{Flujo A:} La interfaz muestra los resultados del análisis seleccionado. \newline \textbf{Flujo B:} La lista del historial de análisis queda vacía. \\
		\textbf{Excepciones}          & Ninguna. \\
		\textbf{Importancia}          & Media \\
		\bottomrule
	\end{tabularx}
	\caption{CU-3 Gestionar el historial de análisis.}
\end{table}

\imagen{use_cases_diagram}{Diagrama de casos de uso de la aplicación de análisis de \textit{malware}.}{1}

\apendice{Especificación de diseño}

\section{Introducción}

\section{Diseño de datos}
% Usar https://www.appdome.com/how-to/devsecops-automation-mobile-cicd/appdome-basics/structure-of-an-android-app-binary-apk como referencia
% Explicación de la estructura de una APK. Explicar que es una especie de archivo comprimido basado en un jar. 
% Usar una estructura similar y explicar cada una de las diferentes partes que puede tener,
\dirtree{%
	.1 /.
	.2 assets.
	.2 lib.
	.3 arm64-v8a.
	.3 x86.
	.3 x86\_64.
	.3 \dots{}.
	.2  META-INF.
	.3 CERT.SF (might have a different name).
	.3 CERT.RSA (might have a different name).
	.3 MANIFEST.MF.
	.2 kotlin.
	.2 res.
	.2 AndroidManifest.xml.
	.2 classes.dex.
	.2 resources.arsc.
}

% Diseño de los datasets usados. Drebin y el dataset propio. Explicar en cada uno de ellos su formato y de dónde viene cada una de esas categorías. Basarse en el esquema y explicación anterior para no tener que detallar mucho el proceso de extracción. En esencia, para esta sección no es necesario entrar en detalle de cómo se extrae o realiza el proceso. Es únicamente para detallar estructuras de datos y diseño que se han usado.

% Estructura de Drebin antes de tranformar:
\dirtree{%
	.1 drebin.
	.2 apks.
	.3 drebin\_0.zip.
	.3 drebin\_1.zip.
	.3 \dots{}.
	.2 metadata.
	.3 dataset\_splits.
	.3 drebin\_results.
	.3 feature\_vectors.
	.4 000a067df9235ae\dots{}.
	.4 \dots{}.
	.3 drebin.libsvm.
	.3 sha256\_family.csv.
}

% Estructura de Drebin tras transformar (mostrar fila de ejemplo de alguna forma):
% sha256 | hw_features | req_permissions | app_components | intents | api_calls | used_permissions | urls | malware | family
% 0002ab | android.hardware.tourchscreen, ... | android.permissions-READ_PHONE_STATE, ... | .GameBootReceiver,.GameService, ... | android.intent.action.BOOT_COMPLETED,android.intent.category.HOME, ... | android/content/Context;->startService,android/app/Activity;->startActivity, ... | android.permission.INTERNET,android.permission.VIBRATE, ... | Read/Write External Storage,getPackageInfo, ... | http://m-001.net/i/,m-001.net, ... | 1 | GinMaster

% Estructura del dataset propio (mosrtar fila de ejemplo de alguna forma):
% 	file_size	fuzzy_hash	activities_list	services_list	receivers_list	permissions_list	api_calls_list	opcode_counts	is_malware
% 34234 | asdasdb | [android..., ...] | [] | ...

% Explicar pipeline de preprocesado (externo) tokenizado, padding, escalado (interno) embedding

\section{Diseño procedimental}
% Diferentes diagramas de secuencia para la aplicación, explicar casos de uso principalmetne. Usar código mermaid para generar los diagramas.

\section{Diseño arquitectónico}

% diseño del modelo (capas) explicarlo
%APKAnalysisModel( 
%	(embedder): APKFeatureEmbedder( # Embedding dim configurable tambien; todos las caracteristicas usan la misma embedding dim
%		(seq_embedders): ModuleDict( 
%			(activities_list): Embedding(247562, 64, padding_idx=0)
%			(services_list): Embedding(22045, 64, padding_idx=0)
%			(receivers_list): Embedding(15922, 64, padding_idx=0)
%			(permissions_list): Embedding(7781, 64, padding_idx=0)
%			(api_calls_list): Embedding(3843329, 64, padding_idx=0)
%		)
%		(char_embedders): ModuleDict(
%			(fuzzy_hash): Embedding(67, 32, padding_idx=0)
%		)
%		(char_gru): ModuleDict(
%			(fuzzy_hash): GRU(32, 64, batch_first=True)
%		)
%		(vector_reducers): ModuleDict(
%			(opcode_counts): Sequential(
%				(0): Linear(in_features=768, out_features=128, bias=True)
%				(1): ReLU()
%				(2): Linear(in_features=128, out_features=64, bias=True)
%				(3): ReLU()
%			)
%		)
%	)
%	(classifier): APKClassifier( 
%		(mlp): Sequential(
%			(0): Linear(in_features=449, out_features=64, bias=True)
%			(1): ReLU()
%			(2): Dropout(p=0.5, inplace=False)
%			(3): Linear(in_features=64, out_features=2, bias=True) # Dinámico en función de un hiperparámetro hidden_dims. Tanto número de capas cópmo neurons
%		)
%	)
%)

% Diagrama del modelo

% diseño de la app de streamlit, modelo cliente servidor


%%%

\apendice{C}{Especificación de Diseño}
\label{apendice:diseno}

\section{Introducción}
Este apéndice presenta el diseño técnico detallado del sistema desarrollado. Su objetivo es servir como un plano que describe la arquitectura del software, la estructura de los datos con los que opera y la lógica procedimental de sus operaciones clave. Se abordará el diseño desde tres perspectivas: el diseño de datos, que define la estructura de los archivos APK y los \textit{datasets}; el diseño procedimental, que ilustra las interacciones principales mediante diagramas de secuencia; y el diseño arquitectónico, que desglosa la estructura interna del modelo de inteligencia artificial y de la aplicación web.

%

Este apéndice detalla el diseño técnico del sistema desarrollado, abarcando tres áreas fundamentales: el diseño de los datos, el diseño procedimental y el diseño arquitectónico. El objetivo es proporcionar una descripción clara y estructurada de cómo se ha organizado la información, cómo fluyen los procesos lógicos y cómo se ha construido la arquitectura del software, desde el modelo de inteligencia artificial hasta la aplicación web. Este documento sirve como un plano técnico que facilita la comprensión, el mantenimiento y la futura expansión del proyecto.

\section{Diseño de datos}
El diseño de datos es fundamental para cualquier sistema de aprendizaje automático. En esta sección se describe la estructura de los datos de entrada primarios (archivos APK) y los conjuntos de datos estructurados que se han derivado de ellos para el entrenamiento y la evaluación de los modelos.

%

El diseño de datos es fundamental para cualquier sistema de aprendizaje automático. En esta sección se describe la estructura de los datos de entrada (archivos APK), el formato de los \textit{datasets} utilizados durante el desarrollo y el \textit{pipeline} de preprocesamiento que los transforma para que puedan ser consumidos por los modelos.

\subsection{Estructura de un archivo APK}
Un archivo APK (\textit{Android Package Kit}) es el formato de paquete utilizado por el sistema operativo Android para la distribución e instalación de aplicaciones móviles. Aunque parece un único archivo, en realidad es un archivo comprimido (basado en el formato JAR y, por tanto, compatible con ZIP) que contiene un conjunto de ficheros y directorios con una estructura bien definida. La estructura de una APK es la siguiente:\\

\dirtree{%
	.1 /.
	.2 assets.
	.2 lib.
	.3 arm64-v8a.
	.3 x86.
	.3 x86\_64.
	.3 \dots{}.
	.2 META-INF.
	.3 CERT.SF.
	.3 CERT.RSA.
	.3 MANIFEST.MF.
	.2 kotlin.
	.2 res.
	.2 AndroidManifest.xml.
	.2 classes.dex.
	.2 resources.arsc.
}

Los componentes más relevantes son:
\begin{itemize}
	\item \textbf{\texttt{assets/}}: Contiene recursos brutos que la aplicación puede utilizar, como ficheros de configuración, bases de datos o recursos de \textit{machine learning}.
	
	\item \textbf{\texttt{lib/}}: Como Android es multiplataforma, este directorio incluye el código compilado específico de las diferentes arquitecturas de procesador (ARM, \texttt{x86}). Contiene, a su vez, las librerías nativas de la aplicación.

	\item \textbf{\texttt{META-INF/}}: Almacena los metadatos de la firma de la aplicación. Los ficheros \texttt{CERT.SF} y \texttt{CERT.RSA} contienen el certificado y la firma que garantizan la autenticidad e integridad de la aplicación, mientras que \texttt{MANIFEST.MF} contiene los \textit{hashes} de todos los ficheros del paquete.	

	\item \textbf{kotlin:} Directorio que contiene el código fuente de Kotlin, si la aplicación está escrita en este lenguaje.
	
	\item \textbf{\texttt{res/}}: Contiene los recursos de la aplicación que no están compilados, como los diseños de la interfaz (\textit{layouts}), las imágenes (\textit{drawables}) o las cadenas de texto (\textit{strings}).
	
	\item \textbf{\texttt{AndroidManifest.xml}:} Es el archivo más importante del APK. Es un fichero XML obligatorio que describe la información esencial sobre la aplicación al sistema Android, como su nombre, componentes (actividades, servicios), permisos, versiones SDK mínima y objetivo.
	
	\item \textbf{\texttt{classes.dex}:} Contiene el código de la aplicación compilado en formato DEX (\textit{Dalvik Executable}), que es el que ejecuta la máquina virtual de Android (ART / Dalvik). Puede haber múltiples archivos .dex si la aplicación es grande.
	
	\item \textbf{\texttt{resources.arsc}:} Es un archivo que contiene recursos precompilados, como las cadenas de texto, para un acceso más eficiente por parte del sistema.
\end{itemize}

\subsection{Diseño de los conjuntos de datos}
A lo largo del proyecto se han utilizado dos \textit{datasets} principales: el \textit{dataset} público Drebin para la fase de prototipado y un \textit{dataset} propio para el desarrollo del modelo final.

\subsubsection{Dataset Drebin}
El formato original del \textit{dataset} Drebin no era una tabla, sino una estructura de directorios compleja. Tras un proceso de \textit{parsing}, se transformó en un único archivo CSV con la siguiente estructura por fila:

\begin{verbatim}
	sha256 | req_permissions | app_components | ... | malware
	000a06 | ["INTERNET", ...] | [".GameService", ...] | ... | 1
\end{verbatim}

Cada fila representa una aplicación, identificada por su \textit{hash} SHA256, y cada columna contiene una lista de las características estáticas extraídas de la misma.

\subsubsection{Dataset propio}
El \textit{dataset} final se construyó desde cero extrayendo características de 20.000 APKs obtenidas de AndroZoo. El formato resultante es un único fichero CSV donde cada fila corresponde a una aplicación. La estructura es la siguiente:
\begin{verbatim}
	file_size | fuzzy_hash     | activities_list   | ... | opcode_counts         | is_malware
	34234     | "ab34...cd56"  | ["com.app...", ...] | ... | [221, 3455, 4567, ...] | 0
\end{verbatim}

\subsection{Pipeline de preprocesado de datos}
Para que los datos de los \textit{datasets} puedan ser utilizados por los modelos, necesitan pasar primero por un \textit{pipeline} de preprocesamiento dividido en dos fases:

\begin{enumerate}
	\item \textbf{Procesamiento externo (offline):} Antes de iniciar cualquier entrenamiento, se realizan una serie de pasos sobre el \textit{dataset} completo. Para cada característica categórica (como los permisos o las actividades), se construye un \textbf{vocabulario} que asigna un índice numérico único a cada posible valor. Luego, todas las listas de cadenas de texto se convierten en listas de índices. Finalmente, se aplica un \textbf{padding} para que todas las listas de una misma columna tengan la misma longitud, convirtiéndolas en matrices numéricas.	
	
	\item \textbf{Procesamiento interno (online):} Una vez los datos están en formato de matrices de índices, el \textit{embedder} del modelo se encarga del resto. Durante el entrenamiento o la inferencia, aplica las capas de \textit{embedding} para convertir estos índices en vectores densos. Además, las características numéricas escalares (como el \texttt{file\_size}) son normalizadas internamente por el modelo para que sus rangos de valores no desestabilicen el entrenamiento.
\end{enumerate}

\section{Diseño procedimental}
Para ilustrar el flujo de trabajo y las interacciones dentro del sistema, a continuación se presentan los diagramas de secuencia de los dos procesos más importantes: el análisis de un APK por parte de un usuario y el proceso interno de entrenamiento de un modelo.

\subsubsection{Diagrama de secuencia: Análisis de un nuevo archivo APK}
Este diagrama muestra la secuencia de interacciones desde que un usuario sube un archivo a la aplicación web hasta que recibe una predicción.

%	sequenceDiagram
%	actor Usuario
%	participant AppWeb as Aplicación Web (Streamlit)
%	participant Backend as Backend (Pipeline de IA)
%	
%	Usuario->>+AppWeb: Sube un archivo APK
%	Usuario->>+AppWeb: Presiona el botón "Analizar"
%	AppWeb->>+Backend: solicitaAnalisis(archivoAPK)
%	Backend->>Backend: extraerCaracteristicas()
%	Backend->>Backend: preprocesarDatos()
%	Backend->>Backend: realizarInferencia()
%	Backend-->>-AppWeb: devuelvePrediccion(resultado)
%	AppWeb-->>-Usuario: Muestra los resultados en la interfaz

\imagen{sequence_predict}{Diagrama de secuencia del proceso de analizar una APK.}

\subsubsection{Diagrama de secuencia: Proceso de entrenamiento del modelo}
Este diagrama detalla el bucle de entrenamiento de la red neuronal, un proceso iniciado por el desarrollador para ajustar los pesos del modelo.

%	sequenceDiagram
%	actor Desarrollador
%	participant Script as Script de Entrenamiento
%	participant DataLoader as Cargador de Datos
%	participant Modelo as Red Neuronal (Embedder + Clasificador)
%	participant Optimizador as Optimizador (Adam)
%	
%	Desarrollador->>+Script: Iniciar entrenamiento()
%	loop Para cada Época
%	Script->>+DataLoader: Solicitar lote de datos
%	DataLoader-->>-Script: Entrega lote (datos, etiquetas)
%	Script->>+Modelo: forward(datos)
%	Modelo-->>-Script: Devuelve predicciones
%	Script->>Script: calcularPerdida(predicciones, etiquetas)
%	Script->>+Modelo: backward()
%	Modelo-->>-Script: Calcula gradientes
%	Script->>+Optimizador: step()
%	Optimizador-->>-Modelo: Actualiza los pesos
%	Script->>-DataLoader: (fin de lote)
%	end
%	Script-->>-Desarrollador: Entrenamiento finalizado

\imagen{sequence_training}{Diagrama de secuencia del proceso de entrenamiento del modelo de red neuronal.}

\section{Diseño arquitectónico}
En esta sección se detalla la arquitectura tanto del modelo de inteligencia artificial como de la aplicación web de demostración.

\subsection{Arquitectura del modelo}
En esta sección se detalla la arquitectura del software, tanto del modelo de inteligencia artificial como de la aplicación web que lo utiliza.

\subsection*{Arquitectura del modelo}
El modelo de inteligencia artificial, implementado en PyTorch, sigue una arquitectura modular y flexible, compuesta por dos grandes bloques: el \textit{embedder} y el clasificador.

\begin{verbatim}
	APKAnalysisModel(
			(embedder): APKFeatureEmbedder(
			(seq_embedders): ModuleDict(...)
			(char_embedders): ModuleDict(...)
			(char_gru): ModuleDict(...)
			(vector_reducers): ModuleDict(...)
		)
		(classifier): APKClassifier(
			(mlp): Sequential(...)
		)
	)
\end{verbatim}

\begin{itemize}
	\item \textbf{APKFeatureEmbedder (Embedder):} Es el componente más complejo y el corazón del sistema. Su única responsabilidad es recibir las características preprocesadas de un APK y convertirlas en un único vector numérico denso. Para ello, contiene diferentes submódulos especializados:
	\begin{itemize}
		\item \textbf{seq\_embedders:} Un diccionario de capas `Embedding` de PyTorch, una para cada característica de tipo lista (permisos, actividades, etc.). Cada capa aprende a representar los elementos de su vocabulario como un vector.
		\item \textbf{char\_embedders y char\_gru:} Un módulo especializado para procesar el \texttt{FUZZY\_HASH}. Trata el \textit{hash} como una secuencia de caracteres, los convierte en vectores con un \textit{embedding} y luego los procesa con una capa GRU (Gated Recurrent Unit) para capturar patrones secuenciales.
		\item \textbf{vector\_reducers:} Un diccionario de pequeñas redes neuronales (MLPs) que toman las características que ya son vectores numéricos (como \texttt{OPCODE\_COUNTS}) y reducen su dimensionalidad para que sea consistente con la de los otros \textit{embeddings}.
	\end{itemize}
	La salida de todos estos submódulos, junto con las características escalares, se concatena para formar el vector final.
	
	\item \textbf{APKClassifier (Clasificador):} Este componente es la cabeza del modelo. Por defecto, es una red neuronal de tipo Perceptrón Multicapa (MLP) que toma el vector del \textit{embedder} y, a través de una o más capas ocultas con funciones de activación ReLU y capas de \textit{dropout}, lo procesa para obtener la predicción final en sus dos neuronas de salida.
\end{itemize}

Una de las claves de este diseño es que la cabeza clasificadora es **intercambiable**. La salida del \textit{embedder} es un vector de características de alta calidad que puede ser utilizado para entrenar cualquier otro modelo de aprendizaje automático clásico (como RandomForest, XGBoost, SVM, etc.), tratando a la red neuronal simplemente como un potente paso de ingeniería de características.

%\imagen{sequence_model_arch}{Diagrama de la arquitectura modelo de red neruonal.}

\subsection{Arquitectura de la aplicación web}

La aplicación de demostración se ha desarrollado con Streamlit y sigue una arquitectura cliente-servidor simple.

\begin{itemize}
	\item \textbf{Cliente:} Es el navegador web del usuario. Se encarga de renderizar la interfaz de usuario y de enviar las interacciones del usuario (como la subida de un archivo o el clic en un botón) al servidor.
	
	\item \textbf{Servidor:} Es el \textit{script} de Python de Streamlit que se ejecuta en el servidor. Este se encarga de toda la lógica de la aplicación: recibe las peticiones del cliente, carga los modelos de IA de los artefactos guardados, ejecuta el \textit{pipeline} de análisis sobre los datos recibidos y genera dinámicamente los componentes de la interfaz (tablas, gráficos, texto) que se envían de vuelta al cliente para ser mostrados.
\end{itemize}

%%

\subsection*{Arquitectura de la aplicación web}
La aplicación de demostración se ha desarrollado con Streamlit y sigue una arquitectura cliente-servidor simple.
\begin{itemize}
	\item \textbf{Cliente (Navegador Web):} El usuario interactúa con la interfaz de la aplicación, que se renderiza en su navegador. Esta interfaz está compuesta por los componentes que ofrece Streamlit (botones, selectores, pestañas, etc.).
	\item \textbf{Servidor (Script de Python):} Un único \textit{script} de Python se ejecuta en el servidor. Este \textit{script} contiene toda la lógica de la aplicación. Se encarga de cargar los modelos entrenados en memoria al iniciarse, gestionar las interacciones del usuario (como la subida de archivos), llamar al \textit{pipeline} de análisis para obtener las predicciones y generar los gráficos de interpretabilidad. Cuando el cliente realiza una acción, el servidor la procesa y le devuelve los nuevos elementos visuales para que la interfaz se actualice.
\end{itemize}
Este diseño permite un desarrollo muy rápido y una integración perfecta con todo el código de análisis y modelado desarrollado en Python.

\apendice{Documentación técnica de programación}
\label{apendice:documentacion_tecnica}

\section{Introducción}
Este apéndice sirve como manual técnico y guía para cualquier desarrollador que desee comprender, extender o contribuir a este proyecto. El objetivo es proporcionar toda la información necesaria para configurar el entorno de desarrollo, entender la estructura del código fuente, y ejecutar los procesos clave como la creación del \textit{dataset} o el entrenamiento de los modelos. Se asume que el lector tiene conocimientos de Python y está familiarizado con herramientas de desarrollo como Git y los entornos virtuales.

\section{Estructura de directorios}
A continuación se muestra y explica la estructura de directorios del repositorio principal del proyecto, alojado en GitHub\footnote{\url{https://github.com/dtx1007/TFG-Malware-Detection-Android}}.

\dirtree{%
	.1 /.
	.2 apks/\\
	{	
		\hspace*{1em}%
		\begin{minipage}[t]{0.85\textwidth}
			\normalfont
			Directorio destinado a contener los archivos APK en bruto descargados para la creación del \textit{dataset}. También almacena el catálogo de AndroZoo (\texttt{latest.csv}). Por su gran tamaño, el contenido de este directorio no se incluye en el repositorio de Git.
		\end{minipage}
	}.
}

\newpage

\dirtree{%
	.1 /.
	.2 dataset/\\
	{	
		\vspace*{0.5em}%
		\hspace*{1em}%
		\begin{minipage}[t]{0.85\textwidth}
			\normalfont
			Directorio donde se colocarían los \textit{datasets} procesados en formato \texttt{.csv} o \texttt{.pkl}, para ser utilizados por los cuadernos de entrenamiento. No se incluyen dentro del repositorio por motivos de tamaño pero se explica como generar el \textit{dataset} propio más adelante.
		\end{minipage}
	}.
	.2 doc/\\
	{	
		\vspace*{0.5em}%
		\hspace*{1em}%
		\begin{minipage}[t]{0.85\textwidth}
			\normalfont
			Alberga toda la documentación del proyecto, incluyendo la memoria y los anexos en formato \LaTeX.
		\end{minipage}
	}.
	.2 model\_artifacts/\\
	{	
		\vspace*{0.5em}%
		\hspace*{1em}%
		\begin{minipage}[t]{0.85\textwidth}
			\normalfont
			Directorio donde se guardan todos los artefactos generados tras el entrenamiento: los modelos de red neuronal y clásicos, los vocabularios de características, los <<escaladores>> de normalización y las métricas de rendimiento.
		\end{minipage}
	}.
	.2 plots/\\
	{	
		\vspace*{0.5em}%
		\hspace*{1em}%
		\begin{minipage}[t]{0.85\textwidth}
			\normalfont
			Contiene todas las gráficas generadas para el análisis de resultados y la interpretabilidad de los modelos.
		\end{minipage}
	}.
	.2 src/\\
	{	
		\vspace*{0.5em}%
		\hspace*{1em}%
		\begin{minipage}[t]{0.85\textwidth}
			\normalfont
			El corazón del proyecto. Contiene todo el código fuente en Python y los \textit{notebooks} de pruebas de Jupyter.
		\end{minipage}
	}.
}

\section{Manual del programador}
Esta sección está pensada como una guía de inicio rápido para un desarrollador. Se detallará la estructura del código fuente y se explicarán los flujos de trabajo más comunes, como la creación del \textit{dataset} y el entrenamiento de los modelos.

Es importante destacar que este repositorio no incluye los \textit{datasets} en bruto (APKs) ni los modelos ya entrenados debido a su gran tamaño. Para una demostración rápida con los modelos ya listos, se recomienda utilizar el \href{https://github.com/dtx1007/streamlit_malware_detection_app}{repositorio de despliegue}. Esta guía está orientada a quien desee replicar el proceso desde cero.

\subsection{Explicación del código fuente}
Todo el código se encuentra en el directorio \texttt{src/}, organizado en subdirectorios con responsabilidades claras.

\subsubsection{Directorio \texttt{src/app}}
Contiene todo el código de la aplicación web interactiva desarrollada con Streamlit.

\dirtree{%
	.1 app/.
	.2 malware\_interpreter\_app.py.
	.2 model\_utils.py.
	.2 data\_utils.py.
	.2 shap\_utils.py.
	.2 ui\_utils.py.
}

\begin{itemize}
	\item \textbf{\texttt{malware\_interpreter\_app.py}:} Es el punto de entrada de la aplicación web. Inicializa y controla la interfaz de usuario, gestiona el estado de la sesión (como el historial) y coordina las llamadas al resto de módulos para realizar el análisis cuando un usuario sube un APK.
	
	\item \textbf{\texttt{model\_utils.py}:} Encapsula la lógica de carga y uso de los modelos de IA. Contiene funciones para cargar los modelos (tanto para la red neuronal como para los modelos clásicos) desde disco usando un sistema de caché para optimizar el rendimiento, y para ejecutar las predicciones.
	
	\item \textbf{\texttt{data\_utils.py}:} Gestiona la carga de datos auxiliares necesarios para la aplicación, como el \textit{dataset} de fondo que utiliza SHAP para sus explicaciones.
	
	\item \textbf{\texttt{shap\_utils.py}:} Centraliza la lógica de interpretabilidad. Contiene las funciones para crear el <<explicador>> de SHAP y para calcular los valores de importancia de las características para una predicción concreta.
	
	\item \textbf{\texttt{ui\_utils.py}:} Módulo con funciones de ayuda para generar componentes visuales específicos, como los gráficos de SHAP o la proyección UMAP.
\end{itemize}

\subsubsection{Directorio \texttt{src/notebooks}}
Contiene los Jupyter Notebooks, que sirven a modo de <<campo de pruebas>> para el proyecto, se utilizan para la experimentación, el prototipado y el análisis de datos.

\newpage

\dirtree{%
	.1 notebooks/.
	.2 download\_apks.ipynb.
	.2 androguard\_feature\_extraction.ipynb.
	.2 models\_training.ipynb.
	.2 model\_intrepretability.ipynb.
	.2 models\_stats.ipynb.
	.2 quantize\_model.ipynb.
	.2 hyperparameters\_optimization.ipynb.
}

\begin{itemize}
	\item \textbf{\texttt{download\_apks.ipynb}:} Contiene el código para descargar los archivos APK desde AndroZoo, utilizando su catálogo para filtrar y seleccionar las muestras.
	
	\item \textbf{\texttt{androguard\_feature\_extraction.ipynb}:} Es uno de los cuadernos más importantes. Implementa el \textit{pipeline} que toma los APKs descargados, los analiza con Androguard para extraer sus características estáticas y crea con ellas el \textit{dataset} final.
	
	\item \textbf{\texttt{models\_training.ipynb}:} Contiene toda la lógica para entrenar y evaluar tanto la red neuronal como los modelos clásicos, utilizando la validación cruzada. Al finalizar, guarda los modelos y artefactos en disco.
	
	\item \textbf{\texttt{model\_intrepretability.ipynb}:} Se centra en la interpretabilidad. Carga los modelos ya entrenados y utiliza SHAP para generar los datos necesarios para las visualizaciones de diferentes gráficos como el de importancia de características, gráficos de fuerza, gráficos en cascada, etc.
	
	\item \textbf{\texttt{models\_stats.ipynb}:} Permite obtener una vista general del rendimiento de los modelos entrenados. Carga las métricas de todos los modelos entrenados y genera varios gráficos comparativos (boxplots, curvas ROC/PR, etc.).
	
	\item \textbf{\texttt{quantize\_model.ipynb}:} Cuaderno de pruebas para aplicar técnicas de cuantización al modelo de red neuronal y evaluar su impacto en el tamaño y el rendimiento.
	
	\item \textbf{\texttt{hyperparameters\_optimization.ipynb}:} Implementa el proceso de búsqueda de hiperparámetros óptimos utilizando la librería Optuna.
\end{itemize}

\subsubsection{Directorio \texttt{src/prototypes}}
Este es el directorio más importante a nivel de código fuente. Contiene los módulos de Python que definen la arquitectura y la lógica de los modelos de IA.
\dirtree{%
	.1 prototypes/.
	.2 torch\_apk\_analysis\_model.py.
	.2 torch\_apk\_analysis\_model\_io.py.
	.2 ml\_model.py.
	.2 ml\_model\_io.py.
}
\begin{itemize}
	\item \textbf{\texttt{torch\_apk\_analysis\_model.py}:} Define la arquitectura completa de la red neuronal en PyTorch, incluyendo las clases para el \texttt{APKFeatureEmbedder}, el \texttt{APKClassifier} y el modelo principal que los une. También contiene las funciones para el bucle de entrenamiento, la evaluación y la inferencia.
	
	\item \textbf{\texttt{torch\_apk\_analysis\_model\_io.py}:} Gestiona el guardado y la carga del modelo de red neuronal y todos sus metadatos asociados (vocabularios, <<escaladores>>, etc.).
	
	\item \textbf{\texttt{ml\_model.py}:} Define y entrena los modelos de aprendizaje automático clásicos (RandomForest, XGBoost, etc.) utilizando la librería scikit-learn.
	
	\item \textbf{\texttt{ml\_model\_io.py}:} Se encarga de guardar y cargar los modelos clásicos ya entrenados.
\end{itemize}

\subsubsection{Directorio \texttt{src/utils}}
Contiene \textit{scripts} con funciones de utilidad que dan soporte al resto del proyecto.
\dirtree{%
	.1 utils/.
	.2 feature\_extraction.py.
	.2 preprocessing\_utils.py.
	.3 vocab\_utils.py.
}
\begin{itemize}
	\item \textbf{\texttt{feature\_extraction.py}:} Contiene la lógica del proceso de extracción de características con Androguard. Es utilizado por la aplicación web para analizar nuevas APKs.
	
	\item \textbf{\texttt{preprocessing\_utils.py}:} Contiene las funciones para preprocesar los datos extraídos (creación de vocabularios, \textit{padding}, etc.), asegurando que el tratamiento de los datos sea consistente en todo el proyecto.
	
	\item  \textbf\texttt{{vocab\_utils.py}:} Contiene funciones para la generación de los vocabularios empleados en el preprocesamiento.
\end{itemize}

\subsection{Guía de procesos comunes} \label{procesos_comunes}
A continuación se describen los pasos para ejecutar las tareas más comunes. Se asume que el entorno de desarrollo ya está configurado (ver sección \ref{entorno_desarrollo}).

\subsubsection{Proceso de creación del \textit{dataset}}
\begin{enumerate}
	\item \textbf{Obtener clave de API:} Es necesario solicitar una clave de API de AndroZoo\footnote{Instrucciones de acceso disponibles en: \url{https://androzoo.uni.lu/access}}. Una vez obtenida, crear un archivo \texttt{.env} en la raíz del proyecto con el formato \texttt{ANDROZOO\_API\_KEY=tu\_clave\_aqui}.
	
	\item \textbf{Descargar catálogo:} Descargar el catálogo actualizado de AndroZoo\footnote{Disponible en: \url{https://androzoo.uni.lu/static/lists/latest.csv.gz}}, descomprimirlo y guardarlo en el directorio \texttt{apks/} con el nombre \texttt{latest.csv}.
	
	\item \textbf{Descargar APKs:} Abrir el cuaderno \\\texttt{src/notebooks/download\_apks.ipynb}. Configurar los parámetros al final del \textit{notebook} (número de muestras, rutas de guardado) y ejecutarlo. El proceso comenzará a descargar las aplicaciones.
	
	\item \textbf{Extraer características:} Abrir el cuaderno \\\texttt{src/notebooks/androguard\_feature\_extraction.ipynb}. Asegurarse de que las rutas de entrada coinciden con las del paso anterior y ejecutar la celda correspondiente para crear el \textit{dataset}. El resultado se guardará por defecto en el directorio \texttt{dataset/}.
\end{enumerate}

\subsubsection{Proceso de entrenamiento del modelo}
Con el \textit{dataset} ya creado, el entrenamiento se realiza desde el cuaderno \texttt{src/notebooks/models\_training.ipynb}. Es importante configurar correctamente las celdas iniciales, especialmente el parámetro \texttt{load\_fresh=True} la primera vez que se use un nuevo \textit{dataset} para que se preprocese y se guarde en disco una versión ya preprocesada. En ejecuciones posteriores, se puede poner en \texttt{False} para ahorrar tiempo. Tras configurar los hiperparámetros deseados, se puede ejecutar el cuaderno completo. Los modelos y artefactos resultantes se guardarán en el directorio \texttt{model\_artifacts/}.

\section{Compilación, instalación y ejecución del proyecto}
En esta sección se detallan los dos flujos principales para ejecutar el proyecto: uno para el desarrollo local y otro para el despliegue mediante Docker.

\subsection{Entorno de desarrollo (local)} \label{entorno_desarrollo}
Este es el flujo recomendado para modificar el código o entrenar nuevos modelos.
\begin{enumerate}
	\item \textbf{Prerrequisitos:} Asegurarse de tener instalados Python 3.11+, Git y Poetry\footnote{Documentación oficial de Poetry: \url{https://python-poetry.org/docs/}}.
	
	\item \textbf{Clonar el repositorio:} Descargar el código fuente desde el repositorio principal.
	
	\item \textbf{Instalar dependencias:} Navegar a la raíz del proyecto en un terminal y ejecutar uno de los siguientes comandos:
	
	\vspace*{-1em}
	\begin{verbatim}
		# Para CPU
		poetry install --extras=cpu
		# Para GPU (con CUDA)
		poetry install --extras=gpu
	\end{verbatim}
	\vspace*{-1em}
	
	Poetry creará automáticamente un entorno virtual e instalará todos los paquetes necesarios. Se puede elegir la versión con o sin soporte para GPU. Es importante aclarar que, se ha de especificar uno de los extras o PyTorch no se instalará, impidiendo que se pueda ejecutar correctamente la aplicación o el resto de \textit{scripts} y \textit{notebooks}.

	\item \textbf{Activar el entorno:} Ejecutar \texttt{poetry shell} para activar el entorno virtual. Una vez dentro, se pueden ejecutar todos los \textit{scripts} y \textit{notebooks} del proyecto.
	
	\item\textbf{ Desplegar aplicación en local (opcional):} Si se quisiera probar la aplicación en local solo haría falta ejecutar el siguiente comando:
	
	\begin{verbatim}
		streamlit run ./src/app/malware_interpreter_app.py
	\end{verbatim}
	
	\vspace{-1em}
	Si justo se ha descargado el repositorio la aplicación dará error al intentar acceder a ella puesto que no tendrá modelos que cargar, se han de entrenar primero (ver \ref{procesos_comunes}).
\end{enumerate}

\subsection{Entorno de despliegue (Docker)}
\label{despliegue_docker}
Este método es el más sencillo para simplemente probar la aplicación final, ya que no requiere instalar dependencias manualmente. Se recomienda usar el <<repositorio de despliegue>>\footnote{\url{https://github.com/dtx1007/streamlit_malware_detection_app}}.

\begin{enumerate}
	\item \textbf{Requisitos previos:} Tener instalado Git, Git LFS y Docker.
	
	\item \textbf{Clonación y construcción:} Clona el repositorio de despliegue y construye la imagen Docker.
	\begin{verbatim}
		# Clona el repositorio y descarga los modelos
		git clone ...
		cd streamlit_malware_detection_app
		git lfs pull
		
		# Construye la imagen (elegir una)
		docker build -t streamlit-malware-app:cpu .
		docker build --build-arg BUILD_TYPE=gpu -t \
		streamlit-malware-app:gpu .
	\end{verbatim}
	
	\item \textbf{Ejecución:} Lanza el contenedor. La aplicación será accesible en \url{http://localhost:8501}.
	\begin{verbatim}
		# Ejecuta el contenedor (elegir una)
		docker run --rm -p 8501:8501 \
		streamlit-malware-app:cpu
		docker run --rm --gpus all -p 8501:8501 \
		streamlit-malware-app:gpu
	\end{verbatim}
\end{enumerate}

\section{Pruebas del sistema}
Es importante señalar que el proyecto no cuenta con una \textit{suite} de pruebas automatizadas (como tests unitarios o de integración con \textit{frameworks} como \texttt{pytest}). La metodología de \textit{testing} se ha basado en un enfoque manual y exploratorio.

Los diferentes Jupyter Notebooks han servido como el principal entorno de pruebas para cada uno de los módulos. Cada cuaderno actúa como una prueba de concepto y de funcionamiento para la lógica que contiene: el cuaderno de extracción de características valida que el análisis con Androguard funciona, el de entrenamiento valida todo el \textit{pipeline} de modelado, y así sucesivamente. Aunque este enfoque no sustituye a un sistema de pruebas automatizado, ha permitido validar cada componente de forma interactiva y asegurar que estos funcionan correctamente y que los resultados son correctos.

\apendice{Documentación de usuario}

\section{Introducción}

\section{Requisitos de usuarios}

\section{Instalación}

\section{Manual del usuario}



\apendice{Anexo de sostenibilización curricular}

%\section{Introducción}
%Este anexo incluirá una reflexión personal del alumnado sobre los aspectos de la sostenibilidad que se abordan en el trabajo.
%Se pueden incluir tantas subsecciones como sean necesarias con la intención de explicar las competencias de sostenibilidad adquiridas durante el alumnado y aplicadas al Trabajo de Fin de Grado.
%
%Más información en el documento de la CRUE \url{https://www.crue.org/wp-content/uploads/2020/02/Directrices_Sosteniblidad_Crue2012.pdf}.
%
%Este anexo tendrá una extensión comprendida entre 600 y 800 palabras.

\section{Introducción}

Siguiendo las <<Directrices para la introducción de la Sostenibilidad en el Curriculum>> de la CRUE~\cite{crue}, este anexo presenta una reflexión personal sobre cómo el presente Trabajo de Fin de Grado se alinea con los principios de un desarrollo sostenible. Aunque a primera vista un proyecto de ciberseguridad y \textit{software} pueda parecer alejado de los conceptos tradicionales de sostenibilidad, un análisis más profundo revela importantes conexiones con sus tres pilares fundamentales: el social, el ambiental y el económico.

El objetivo de este apartado es, por tanto, analizar el impacto del proyecto a través de estas tres dimensiones y reflexionar sobre cómo, a lo largo de su desarrollo, se han adquirido y aplicado de forma práctica algunas de las competencias transversales para la sostenibilidad que la universidad busca fomentar en sus titulados.

\section{Impacto de la sostenibilidad en el proyecto}

La sostenibilidad, entendida de forma holística, busca un equilibrio entre la justicia social, la viabilidad económica y la calidad ambiental. Este proyecto, aunque de naturaleza técnica, tiene implicaciones en cada una de estas áreas.

\subsubsection{Dimensión social y ètica}
Este es, sin duda, el pilar donde el proyecto tiene su impacto más directo y significativo. El desarrollo de una herramienta para la detección de \textit{malware} es, en esencia, un acto de protección social en el ecosistema digital. El \textit{malware} no es solo una molestia técnica; es una herramienta que facilita el robo de datos personales, el fraude financiero, la violación de la privacidad y la desestabilización de las comunicaciones. Al crear un sistema que combate estas amenazas, el proyecto contribuye directamente a la construcción de un entorno digital más seguro y fiable para todos, un aspecto fundamental para la justicia social en la era en la que vivimos. Este enfoque se alinea con el principio ético de la sostenibilidad, que considera la protección y la seguridad de las personas como un objetivo primordial.

\subsubsection{Dimensión ambiental}
La sostenibilidad ambiental en el ámbito del \textit{software} está intrínsecamente ligada a la eficiencia energética y al uso de los recursos computacionales. En este sentido, la decisión de centrar el proyecto en el análisis estático en lugar del dinámico tiene una clara implicación ambiental. El análisis dinámico requiere la ejecución de aplicaciones en entornos controlados (\textit{sandboxes}), que a menudo son máquinas virtuales que consumen una cantidad considerable de energía. El análisis estático, al no requerir ejecución, es un método inherentemente más eficiente y con una menor huella de carbono, especialmente si se piensa en su aplicación a gran escala.

Además, la conclusión final del proyecto, que aboga por un enfoque híbrido donde un \textit{embedder} de red neuronal potencia a clasificadores clásicos mucho más ligeros, también tiene una vertiente de sostenibilidad. Un modelo como la Regresión Logística o XGBoost requiere muchos menos recursos computacionales para realizar una predicción que una red neuronal profunda completa. Esto significa que, si la herramienta se desplegara en millones de dispositivos, el consumo energético agregado sería significativamente menor, contribuyendo a la sostenibilidad del ecosistema digital.

\subsubsection{Dimensión económica}
El impacto económico del \textit{malware} es devastador, causando pérdidas millonarias cada año tanto a empresas como a particulares. Un sistema de detección eficaz como el que se ha desarrollado en este trabajo contribuye a la sostenibilidad económica al prevenir el fraude, proteger los activos digitales y reducir los costes asociados a la recuperación de ciberataques. Asimismo, al tratarse de un proyecto desarrollado con herramientas de código abierto y con la intención de ser compartido bajo una licencia permisiva, se fomenta un modelo de innovación sostenible, donde el conocimiento y las herramientas se comparten para el beneficio de la comunidad, en lugar de permanecer en silos cerrados.

\section{Competencias de sostenibilidad adquiridas y aplicadas}

El desarrollo de este proyecto también ha sido de gran utilidad para adquirir y poner en práctica las competencias transversales para la sostenibilidad definidas por la CRUE.

\begin{itemize}
	\item \textbf{SOS1 - Contextualización crítica del conocimiento:} El proyecto no se limitó a resolver un problema técnico aislado. Fue necesario investigar y comprender la problemática social y económica del \textit{malware} a nivel global, entendiendo cómo una solución tecnológica como esta interactúa con la sociedad y puede contribuir a mitigar un riesgo que afecta a millones de personas.
	
	\item \textbf{SOS2 - Utilización sostenible de recursos:} Esta competencia se aplicó directamente en la elección de la metodología. Se optó por el análisis estático por ser más eficiente en recursos que el dinámico. Además, todo el proceso de optimización del modelo, que culminó en la propuesta de un enfoque híbrido más ligero, fue un ejercicio práctico de búsqueda de la solución más eficiente y, por tanto, más sostenible.
	
	\item \textbf{SOS3 - Participación en procesos comunitarios:} Aunque es un proyecto individual, su desarrollo se ha enmarcado dentro de la comunidad del software de código abierto. Se han utilizado herramientas creadas por la comunidad (Androguard, PyTorch, etc.) y se ha construido sobre el conocimiento de la comunidad investigadora. La intención de licenciar el proyecto de forma abierta es una forma de devolver ese valor a la comunidad, promoviendo la sostenibilidad del conocimiento compartido.
	
	\item \textbf{SOS4 - Aplicación de principios éticos:} La ética ha sido un pilar central. El proyecto se enfoca en proteger a los usuarios, no en explotar datos. El método de análisis estático se eligió, entre otras cosas, porque respeta la privacidad al no requerir la ejecución de la aplicación ni el acceso a los datos personales del usuario. Además, la inclusión de un módulo de interpretabilidad (SHAP) responde a la necesidad ética de que las decisiones de la IA sean transparentes y comprensibles, y no meras "cajas negras".
\end{itemize}


\bibliographystyle{plain}
\bibliography{bibliografiaAnexos}

\end{document}
